\chapter{简单系统的热力学性质}
\label{chap:simple system}

\section{形式关系}

\begin{definition}
	{广延量和强度量}{extensive parameter & intensive parameter}
	将处于同一状态的样品等比例放大，内能$U$、体积$V$、熵$S$和摩尔数$n$也随之等比例增大，这些量称为广延量(extensive parameter)。当子系统间的相互作用和界面贡献可忽略时，熵可加：
	\begin{equation}
		S(U,V,n)=\sum_{i=1}^kS_i(U_i,V_i,n_i).
	\end{equation}
	而温度$T$、压强$p$%和化学势$\mu$
	等物理量在样品等比例放大时保持不变，称为强度量(intensive parameter)。子系统间可自由交换热量时，平衡要求温度相等；可通过无外载荷的活动隔板交换体积时，平衡还要求压强相等。
\end{definition}

\begin{theorem}
	{Euler方程}{Euler's equation}
	设$r$元函数$f$在所讨论的点可微，且对任意$\lambda>0$满足$k$阶齐次条件
	\begin{equation}
		\label{eqn:k-order homogeneous}
		f(\lambda x_1,\lambda x_2,\ldots,\lambda x_r)=\lambda^k f(x_1,x_2,\ldots,x_r),
	\end{equation}
	则应有
	\begin{equation}
		\label{eqn:Euler's equation}
		k f(x_1,x_2,\ldots,x_r)=\sum_{i=1}^r x_i\pv f{x_i}.
	\end{equation}
\end{theorem}

\begin{proof}
	对\eqref{eqn:k-order homogeneous}两边同时对$\lambda$求导
	\[
		x_1\pv f{\lambda x_1}+x_2\pv f{\lambda x_2}+\cdots+x_r\pv f{\lambda x_r}=k\lambda^{k-1}f(x_1,x_2,\ldots,x_r).
	\]
	取$\lambda=1$即得\eqref{eqn:Euler's equation}。
\end{proof}

\begin{corollary}
	对单组分、仅有体积功的均匀系统，忽略表面效应时，$U=U(S,V,n)$对全部广延变量是一阶齐次的，故
	\begin{equation}
		\label{eqn:U=TS-pV}
		U=TS-pV+\mu n.
	\end{equation}
	其中$\mu:=(\p U/\p n)_{S,V}$为摩尔化学势，详见\secref{sec:chemical potential}。即使过程保持$n$不变，Euler方程中仍有$\mu n$。平衡黑体辐射的光子数可自由调节，$\mu=0$，此时$U=TS-pV$。
	定义摩尔内能$u:=U/n$、摩尔熵$s:=S/n$、摩尔体积$v:=V/n$，以及摩尔热容$c_V:=C_V/n$、$c_p:=C_p/n$。
\end{corollary}

\begin{theorem}
	{Gibbs--Duhem关系}{Gibbs-Duhem relation}
	上述单组分系统的强度量满足：
	\begin{equation}
		S\d T-V\d p+n\d\mu=0.
	\end{equation}
\end{theorem}

\begin{proof}
	对\eqref{eqn:U=TS-pV}两边微分，再与$\d U=T\d S-p\d V+\mu\d n$比较即得
	\[
		0=S\d T-V\d p+n\d\mu.
		\qedhere
	\]
\end{proof}

\begin{example}{用形式关系求基本方程}{}
	某热力学系统的内能和状态方程为
	\[
		U=\frac12pV,\qquad T^2=A\frac{U^{3/2}}{Vn^{1/2}}.
	\]
	其中常数$A>0$，求系统的基本方程。

	条件中出现的广延量为$U,V,n$，因此用熵表象，考虑单位摩尔熵$s=s(u,v)$
	\begin{align*}
		\frac1T&=A^{-1/2}u^{-3/4}v^{1/2},\\
		\frac pT&=2A^{-1/2}u^{1/4}v^{-1/2}.
	\end{align*}
	凑微分：
	\begin{gather*}
		\d s=\frac1T\d u+\frac pT\d v=A^{-1/2}(u^{-3/4}v^{1/2}\d u+2u^{1/4}v^{-1/2}\d v)=4A^{-1/2}\d(u^{1/4}v^{1/2}),\\
		\implies
		s=4A^{-1/2}u^{1/4}v^{1/2}+s_0,
	\end{gather*}
	即
	\[
		S=4A^{-1/2}U^{1/4}V^{1/2}n^{1/4}+ns_0.
	\]
\end{example}

\section{热力学其他表象}

在能量表象和熵表象中，广延量是数学上的独立变量，而强度量是被导出的物理量；%这与现实中实验操作的便利性截然不同。
然而，有时强度量更容易测量或控制，
最突出的是熵$S$和温度$T$这一对共轭变量：并不存在可以测量和控制熵的仪器，而温度计和恒温器在实验室中十分常见。
这种情况下更适合将强度量作为独立变量，广延量作为导出量。
这就需要Legendre变换。
% 这样还能够导出很多其它的热力学表象。


\begin{theorem}
	{Legendre变换}{Legendre transformation}
	考虑函数$y=y(x)$，其在$x$处的切线斜率为$p$、纵截距为$\psi$，则存在$(x,y)\mapsto(p,\psi)$：
	\begin{subequations}
		\label{eqn:Legendre transformation}
		\begin{align}
			p&=\dv yx,\\
			\psi&=y-px,
		\end{align}
	\end{subequations}
	称为Legendre变换。
	\begin{notecaptionblock}
		\tikzfigure{legendre-transform}
		\captionof{figure}{Legendre变换}
		\label{fig:Legendre transform}
	\end{notecaptionblock}
\end{theorem}

\begin{corollary}
	在$p=y'$严格单调的区间内，斜率唯一确定切点，从而
	\[
		\d\psi=\d y-(p\d x+x\d p)=-x\d p,
	\]
	可得Legendre逆变换
	\begin{subequations}
		\begin{align}
			x&=-\dv\psi p,\\
			y&=\psi+px.
		\end{align}
	\end{subequations}
\end{corollary}

\begin{corollary}
	推广到多元函数$y(x_1,x_2,\ldots,x_r)$。在梯度映射$\bm x\mapsto\bm p$光滑可逆的区域内，其Legendre变换为
	\[
		\begin{aligned}
			p_i  & =\pv y{x_i}            \\
			\psi & =y-\sum_{i=1}^rp_ix_i.
		\end{aligned}
		\iff
		\begin{aligned}
			x_i & =-\pv\psi{p_i}            \\
			y   & =\psi+\sum_{i=1}^rp_ix_i.
		\end{aligned}
	\]
\end{corollary}

\begin{example}{分析力学中的Legendre变换}{}
	Lagrange量$L(q,\dot q,t)$可以完整刻画一个力学系统，
	其中$q=(q_1,q_2,\ldots,q_r)$为广义坐标，$\dot q=(\dot q_1,\dot q_2,\ldots,\dot q_r)$为广义速度。广义动量定义为
	\[
		p_i:=\pv L{\dot q_i}.
	\]
	如果以动量代替速度作为独立变量，则需要对$\dot q_i$做Legendre变换，引入Hamilton量
	\begin{equation}
		H(q,p,t):=\sum_{i=1}^rp_i\dot q_i-L(q,\dot q,t).
	\end{equation}
	这样$H(q,p,t)$关于广义动量$p_i$的导数便是广义速度
	\[
		\pv H{p_i}=\dot q_i.
	\]
\end{example}

从已有的能量$U=U(S,V)$出发，进行一系列Legendre变换，得到其它热力学表象。
% 从能量$\d U=T\d S-p\d V$出发

\begin{definition}
	{焓}{enthalpy}
	考虑等压过程，$V\to p=-\p U/\p V$，可得焓(enthalpy) $H=H(S,p)$
	\begin{subequations}
		\begin{align}
			H&=U+pV,\\
			\d H&=T\d S+V\d p.
		\end{align}
	\end{subequations}
	% \begin{equation}
	% 	H=U+pV.
	% \end{equation}
\end{definition}

\begin{definition}
	{Helmholtz自由能}{Helmholtz free energy}
	考虑等温过程，$S\to T=\p U/\p S$，可得Helmholtz自由能$F=F(T,V)$
	\begin{subequations}
		\begin{align}
			F&=U-TS,\\
			\d F&=-S\d T-p\d V.
		\end{align}
	\end{subequations}
	% \begin{equation}
	% 	F=U-TS.
	% \end{equation}
\end{definition}

\begin{corollary}
	自发的等温等容过程中$\Delta F=\Delta U-T\Delta S=Q-T\Delta S\leq 0$自由能减少。
\end{corollary}

\begin{definition}
	{Gibbs自由能}{Gibbs free energy}
	考虑等温等压过程，$(S,V)\to(T,p)$，可得Gibbs自由能$G=G(T,p)$
	\begin{subequations}
		\begin{align}
			G&=U-TS+pV,\\
			\d G&=-S\d T+V\d p.
		\end{align}
	\end{subequations}
	% \begin{equation}
	% 	G=U-TS+pV.
	% \end{equation}
\end{definition}

\begin{corollary}
	自发的等温等压过程中$\Delta G=\Delta U-T\Delta S+p\Delta V=Q-T\Delta S\leq 0$自由能减少。
\end{corollary}

\begin{definition}{特性函数}{characterist funtion}
	适当选取自变量，只需一个热力学量就可决定均匀系统的全部热力学性质，这样的函数称为特性函数(characterist funtion)。
	包括$U(S,V),H(S,p),F(T,V),G(T,p)$等。
\end{definition}

\begin{theorem}
	{Maxwell关系}{Maxwell relation}
	由特性函数$U(S,V),F(T,V),H(S,p),G(T,p)$的二阶导可交换：
	\begin{subequations}
		\begin{alignat}{3}
			\label{eqn:Maxwell relation U}
			\pw USV&=\pw UVS&&\implies& -\pvat{p}{S}{V}&=\pvat{T}{V}{S};\\
			\label{eqn:Maxwell relation H}
			\pw HSp&=\pw HpS&&\implies& \pvat{V}{S}{p}&=\pvat{T}{p}{S};\\
			\label{eqn:Maxwell relation F}
			\pw FTV&=\pw FVT&&\implies& \pvat{p}{T}{V}&=\pvat{S}{V}{T};\\
			\label{eqn:Maxwell relation G}
			\pw GTp&=\pw GpT&&\implies& -\pvat{V}{T}{p}&=\pvat{S}{p}{T}.
		\end{alignat}
	\end{subequations}
	称为Maxwell关系。
	可以借助图形记忆，如\figref{fig:Maxwell relation} 所示，
	外围方形四边上的$U,F,H,G$是四个特性函数，每边的两个顶点分别是状态参量，如$U(S,V)$，对角线连接状态参量和导出量，箭头方向表示正负号，如$\d U=T\d S-p\d V$。再由二阶导可交换便可得到Maxwell关系。
	至于\figref{fig:Maxwell relation} 中字母的顺序，则可以通过一句口诀中的首字母记忆。
	\footnote{其实还不如直接记住$\d U=T\d S-p\d V$和$H=U+pV$，剩下的其实很自然得都出来了。}
	\begin{notecaptionblock}
		\tikzfigure{thermodynamic-square}
		\captionof{figure}{\textbf{G}ood \textbf{p}hysicists \textbf{H}ave \textbf{S}tudied \textbf{U}nder \textbf{V}ery \textbf{F}ine \textbf{T}eachers.}
		\label{fig:Maxwell relation}
	\end{notecaptionblock}
\end{theorem}

\begin{corollary}
	考虑Gibbs自由能$G(T,p)$的二阶导：
	\begin{subequations}
		\begin{alignat}{2}
			\pv[2]GT&=-&\pvat{S}{T}{p}&=-\frac{C_p}T,\\
			\pw GTp&=&\pvat{V}{T}{p}&=V\alpha,\\
			\pv[2]Gp&=&\pvat{V}{p}{T}&=-V\kappa_T.
		\end{alignat}
	\end{subequations}
	因此难以直接测量的响应函数(热力学偏微分)都可以用$C_p,\alpha,\kappa_T$来表示。
\end{corollary}

\begin{example}
	{热弹效应}{}
	由Maxwell关系
	\[
		\pvat{S}{p}{T}=-\pvat{V}{T}{p}=-V\alpha.
	\]
	对$\alpha>0$的材料，可逆等温压缩时$\d S=-V\alpha\d p<0$，材料向热库放热；可逆等温膨胀时吸热。
\end{example}

\begin{example}
	{热容差}{}
	由于液体一般很难压缩，等容热容$C_V$相比等压热容$C_p$更难测量，但是我们可以将二者的差用$\alpha,\kappa_T$表示出来。
	继续\exmref{exm:ideal gas U and C} 的推导。
	利用Maxwell关系\eqref{eqn:Maxwell relation F}，
	\[
		C_p-C_V=\biggb{\pvat{U}{V}{T}+p}\pvat{V}{T}{p}=T\pvat{S}{V}{T}\pvat{V}{T}{p}=T\pvat{p}{T}{V}\pvat{V}{T}{p}.
	\]
	代入响应函数可得
	\begin{equation}
		C_p-C_V=-pVT\alpha\beta=\frac{VT\alpha^2}{\kappa_T}.
	\end{equation}
	\tcblower
	类似地，可定义绝热压缩率
	\begin{equation}
		\kappa_S:=-\frac1V\pvat{V}{p}{S},
	\end{equation}
	压缩率的差也存在类似的关系，自证不难
	\begin{equation}
		\kappa_T-\kappa_S=\frac{VT\alpha^2}{C_p}.
	\end{equation}
\end{example}

\begin{example}{}{}
	化简$\pvat{p}{U}{G}=\pvat{U}{p}{G}^{-1}$，由$\d U=T\d S-p\d V$可得
	\[
		\pvat{U}{p}{G}=T\pvat{S}{p}{G}-p\pvat{V}{p}{G},
	\]
	利用\eqref{eqn:XY->XW/YW}和$\d G=-S\d T+V\d p$可得
	\[
		\pvat{U}{p}{G}=-T\fracdisp{\pvat{G}{p}{S}}{\pvat{G}{S}{p}}+p\fracdisp{\pvat{G}{p}{V}}{\pvat{G}{V}{p}}=-T\fracdisp{-S\pvat{T}{p}{S}+V}{-S\pvat{T}{S}{p}+0}+p\fracdisp{-S\pvat{T}{p}{V}+V}{-S\pvat{T}{V}{p}+0}.
	\]
	由响应函数的定义、\eqref{eqn:XY->XW/YW}和Maxwell关系
	\begin{align*}
		\pvat{T}{V}{p}&=\frac1{V\alpha},\\
		\pvat{T}{S}{p}&=\frac{T}{C_p},\\
		\pvat{T}{p}{S}&=-\Biggdivision{\pvat{S}{p}{T}}{\pvat{S}{T}{p}}=\frac{TV\alpha}{C_p},\\
		\pvat{T}{p}{V}&=-\Biggdivision{\pvat{V}{p}{T}}{\pvat{V}{T}{p}}=\frac{\kappa_T}\alpha.
	\end{align*}
	故
	\[
		\pvat{p}{U}{G}=\biggp{-TV\alpha+\frac{VC_p}S+pV\kappa_T-\frac{pV^2\alpha}S}^{-1}.
	\]
\end{example}

\section{最值原理}

热机能输出多少功，取决于它可以和哪些热源交换热量。先考虑封闭系统只与温度为$T_0>0$的热库换热，记$Q$为系统吸收的热量，$W_{\mathrm{out}}$为系统对外输出的总功。

\begin{theorem}
	{最大功原理}{maximum work principle}
	给定初末态和上述热库，系统对外输出的功满足$W_{\mathrm{out}}\leq T_0\Delta S-\Delta U$。可逆过程达到这一上限。
\end{theorem}

\begin{proof}
	由第一定律，$W_{\mathrm{out}}=Q-\Delta U$。热库放出热量$Q$，系统与热库的总熵变为$\Delta S_{\mathrm{tot}}=\Delta S-Q/T_0\geq0$。消去$Q$得
	\[
		W_{\mathrm{out}}=T_0\Delta S-\Delta U-T_0\Delta S_{\mathrm{tot}} .
	\]
	可逆过程的总熵变为零，产生熵会使输出功减少。
\end{proof}

\begin{corollary}
	若系统初末态的温度均为$T_0$，可逆过程输出的总功等于Helmholtz自由能的减少量$-\Delta F$。若还需在环境压强$p_0$下排开体积$\Delta V$，可供其他用途的功须扣去$p_0\Delta V$。系统与给定环境达到平衡时所能输出的最大有用功称为㶲(exergy)。
\end{corollary}

\begin{example}{}{}
	两份各为$n$摩尔、体积固定的理想气体具有相同的常热容$cnR$。它们与外界绝热，通过循环热机交换能量，初温分别为$T_1,T_2$，求热机所能输出的最大功。

	平衡态两系统温度相等$T$，则能量变化
	\[
		\Delta U=cnR(2T-T_1-T_2).
	\]
	为使功最大，应为可逆过程，即
	\[
		\Delta S=cnR\biggp{\ln\frac T{T_1}+\ln\frac T{T_2}}=cnR\ln\biggp{\frac{T^2}{T_1T_2}}=0
		\implies T=\sqrt{T_1T_2}.
	\]
	最大功为
	\[
		W_{\mathrm{out,max}}=-\Delta U=cnR\Bigp{\sqrt{T_1}-\sqrt{T_2}}^2.
	\]
\end{example}
\begin{example}{}{}
	三份体积固定的理想气体具有相同的常热容$cnR$，初温分别为$T_1,T_2,T_3$。用循环热机在它们之间传递能量，整个装置与外界既不换热也不做功，求其中一份气体所能达到的最大温度。

	设目标气体的末温为$T_\hi$。给定$T_\hi$时，另两份气体的总内能固定，因而其温度之和固定。两温相等时乘积最大，其总熵也最大。因此求可达的最高温度时，可以令另两份气体的末温同为$T_\lo$。由能量守恒，
	\[
		T_\hi+2T_\lo=T_1+T_2+T_3=:T_\tot.
	\]
	总熵变
	\[
		\Delta S=cnR\ln\biggp{\frac{T_\lo^2T_\hi}{T_1T_2T_3}}\geq 0
		\implies(T_\tot-T_\hi)^2T_\hi\geq 4T_1T_2T_3.
	\]
	阴影部分便对应满足条件的$T_\hi$的范围(显然，阴影部分总是存在的)，
	进而可得到其最大值$T_{\max}$。
	\begin{notecaptionblock}
		\tikzfigure{hot-reservoir-temperature}
		\captionof{figure}{$T_\hi$满足的条件}
	\end{notecaptionblock}
\end{example}

\begin{theorem}
	{能量最小原理}{minimum energy principle}
	在正温区，固定总熵、体积、物质的量及其他外部约束，平衡态使内能对可变内部参量取最小值。
\end{theorem}

\begin{proof}
	设平衡态的能量和熵为$U_*,S_*$。固定其余外部约束和任意内部参量$X$，有$S(U_*,X)\le S_*$。由于$(\p S/\p U)_X=1/T>0$，在同一正温支上，能量低于$U_*$的态达不到$S_*$。故在给定熵$S_*$的态中，$U_*$是最小能量。

	在严格稳定的平衡点附近，这一原理也可写成微分形式。在给定总内能时，平衡态使熵最大，
	\[
		\pvat{S}{X}{U}=0,\quad\biggp{\pv[2]SX}_U<0.
	\]
	为方便，记能量的一阶导
	\[
		L:=\pvat{U}{X}{S}=-\Biggdivision{\pvat{S}{X}{U}}{\pvat{S}{U}{X}}=-T\pvat{S}{X}{U}=0.
	\]
	再证能量的二阶导为正
	\[
		\biggp{\pv[2]UX}_S=\pvat{L}{X}{S}=\pvat{L}{U}{X}\cancel{\pvat{U}{X}{S}}+\pvat{L}{X}{U}=0-T\biggp{\pv[2]SX}_U>0.
		\qedhere
	\]
\end{proof}

\begin{theorem}
	{Helmholtz自由能最小原理}{}
	与正温热库$T_\rev$接触、总体积和粒子数固定且仅有体积功的系统，移除内部约束后达到平衡态使得Helmholtz自由能最小。
	%平衡态下，系统与热库进行热接触	时，系统的各个无约束的内部参量取值满足在 T = Tr 的条件下最小化 Helmholtz 自由能。
\end{theorem}

\begin{proof}
	在定容、无其他功的条件下，热库放出的热量等于系统增加的内能。由系统和热库的总熵变化，
	\[
		\Delta S_{\rm tot}=\Delta S-\frac{\Delta U}{T_\rev}
		=-\frac{\Delta(U-T_\rev S)}{T_\rev}\geq0 .
	\]
	在总能量给定时，最大化总熵等价于最小化$U-T_\rev S$。与热库达到热平衡的各候选态都有$T=T_\rev$，此量等于系统的Helmholtz自由能$F$。
	
\end{proof}

\begin{corollary}
	平衡态的限制条件为$T=T_\rev$。
	当这个系统被活塞隔开时
	\[
		\d F=\pv F{V_1}\d V_1+\pv F{V_2}\d V_2=(-p_1+p_2)\d V_1=0.
	\]
	因此平衡时$p_1=p_2.$
\end{corollary}

\begin{theorem}
	{焓最小原理}{}
	与恒压库接触、保持熵和粒子数不变且仅有体积功的系统，移除内部约束后达到平衡态使得焓最小。
\end{theorem}

\begin{theorem}
	{Gibbs自由能最小原理}{}
	与恒温恒压库接触、粒子数固定且仅有体积功的系统，移除内部约束后达到平衡态使得Gibbs自由能最小。
\end{theorem}

\section{稳定性}
\label{sec:thermo-stability}

把一个均匀系统分成等大的两部分，固定每部分的物质的量。若从一部分向另一部分转移少量能量和体积，两部分的熵的一阶变化相消。要使均匀态在这些扰动下稳定，总熵的二阶变化应为负。因此在严格稳定的单相区，$S(U,V)$的Hessian矩阵$H_S$负定，即
\begin{equation}
	\pv[2]SU<0,\qquad \det H_S=\pv[2]SU\pv[2]SV-\biggp{\pw SUV}^2>0.
\end{equation}
这两个条件分别约束热容和压缩率。由$S_U=1/T$和$S_V=p/T$，将二次型写成
\[
	\delta^2 S=\delta(1/T)\,\delta U+\delta(p/T)\,\delta V.
\]
利用Maxwell关系$(\p S/\p V)_T=(\p p/\p T)_V$，有
\[
	\delta U=C_V\delta T+\biggb{T\pvat{p}{T}{V}-p}\delta V.
\]
代回后交叉项相消，得到
\begin{equation}
	\delta^2 S=-\frac{C_V}{T^2}(\delta T)^2-\frac{1}{TV\kappa_T}(\delta V)^2.
\end{equation}
在$T>0$时，严格稳定要求$C_V>0$且$\kappa_T>0$。前者使子系统吸收能量后温度升高，后者使等温压缩时压强升高，两种响应都有助于恢复平衡。特别地，$\kappa_T>0$要求等温线的斜率$(\p p/\p V)_T<0$。这是对小扰动的局域判据；判断气液两相中哪一相稳定，还要比较两相的自由能，见\secref{ssec:liquid-gas phase transition}。

\section{样例系统}

\subsection{理想气体}

前面已经给过，在一定温度范围内，
理想气体内能和状态方程为
\begin{subequations}
	\begin{align}
		u&=cRT,\\
		pv&=RT.
	\end{align}
\end{subequations}
% 其中$c$与理想气体自由度有关，$R=\NA\kB=\qty{8.3144}{\joule\per\mole\kelvin}$是普适的理想气体常数。
% \footnote{最新的国际制基本单位\textit{定义}Avogadro常数$\NA\equiv\qty{6.02214076e23}{\per\mole}$；Boltzmann常数$\kB\equiv\qty{1.380649e-23}{\joule\per\kelvin}$，因此理想气体常数也是精确的。}
可以直接得到响应函数
\begin{subequations}
	\begin{align}
		c_V&=\pvat{u}{T}{V}=\frac uT=cR,\\
		c_p&=\pvat{u}{T}{p}+p\pvat{v}{T}{p}=(c+1)R,\\
		\alpha&=\frac1v\pvat{v}{T}{p}=\frac1T,\\
		\beta&=-\frac1p\pvat{p}{T}{v}=-\frac1T,\\
		\kappa_T&=-\frac1v\pvat{v}{p}{T}=\frac1p.
	\end{align}
\end{subequations}
选定参考点$s_0=s(u_0,v_0)$后，摩尔熵$s$的表达式为
\begin{equation}
	s=R\ln\biggb{\biggp{\frac u{u_0}}^c\frac v{v_0}}+s_0.
\end{equation}
摩尔Helmholtz自由能$f=u-Ts$是$T,v$的函数，写成更简洁的形式：
\begin{equation}
	\frac fT=-R\ln\biggb{\biggp{\frac T{T_0}}^c\frac v{v_0}}+\frac{f_0}{T_0}.
\end{equation}
Gibbs自由能$g=u-Ts+pv$是$T,p$的函数：
\begin{equation}
	\frac gT=-R\ln\biggb{\biggp{\frac T{T_0}}^{c+1}\frac{p_0}p}+\frac{g_0}{T_0}.
\end{equation}
% 考虑等温过程，$\Delta G=nRT\ln(p'/p)$，这一简单关系在\chapref{chap:phase}中会被用到。

\subsectionstar{van der Waals气体}

在低温、高压的条件下，理想气体就与实际气体相差较远了。
van der Waals气体是对理想气体的一个比较好的修正：
取参考态$(u_0,v_0)$的摩尔熵为$s_0$。从理想气体的熵式出发，将摩尔内能减去平均吸引势能$-a/v$，并将摩尔体积减去排除体积$b$，得到
\begin{equation}
	s=s_0+R\ln\biggb{\biggp{\frac{u+a/v}{u_0+a/v_0}}^c\frac{v-b}{v_0-b}} .
\end{equation}
由此可得到van der Waals气体的内能和状态方程：
\begin{subequations}
	\begin{gather}
		u=cRT-\frac av,\\
		\Bigp{p+\frac a{v^2}}(v-b)=RT.
	\end{gather}
\end{subequations}
在力学稳定的均匀分支上，进而得到响应函数
\begin{subequations}
	\begin{align}
		c_V&=cR,\\
		c_p&=\biggp{c+\frac{RTv^3}{RTv^3-2a(v-b)^2}}R>(c+1)R,\\
		\alpha&=\frac{Rv^2(v-b)}{RTv^3-2a(v-b)^2},\\
		\kappa_T&=\frac{v^2(v-b)^2}{RTv^3-2a(v-b)^2}.
	\end{align}
\end{subequations}
在低密度极限$v\to\infty$下，形式变成$pv=RT$，$\alpha=1/T$，$\kappa_T=v/RT=1/p$，变成理想气体情形。

\subsectionstar{黑体辐射}

在一个黑体辐射系统中，内能和状态方程为
\begin{subequations}
	\begin{align}
		U&=\frac{4\sigma}cVT^4,\\
		p&=\frac13\frac UV=\frac{4\sigma}{3c}T^4.
	\end{align}
\end{subequations}
其中$\sigma$是Stefan--Boltzmann常数，$c$在这里是光速。

可得相应函数(恒压过程同时也为恒温过程，$C_p,\alpha,\kappa_T$无意义)：
\begin{subequations}
	\begin{align}
		C_V&=\pvat{U}{T}{V}=4\frac UT=\frac{16\sigma}cVT^3,\\
		\beta&=-\frac1p\pvat{p}{T}{V}=-\frac4T.
	\end{align}
\end{subequations}
熵的微分$\d S=(\d U+p\d V)/T$，积分可得熵
\begin{equation}
	S=\frac43\frac UT=\frac{16\sigma}{3c}VT^3.
\end{equation}
因此可逆绝热过程满足$VT^3=\const$；
Gibbs自由能
\begin{equation}
	G=U-TS+pV=0.
\end{equation}

\subsection{电介质和磁介质}

\paragraph{电介质}

设电场方向固定，$p_\electric$为样品总电偶极矩沿电场的分量。以场强$E$为自变量，Gibbs自由能满足
\[
	\d G=-S\d T+V\d p-p_\electric\d E
\]
Maxwell关系
\begin{equation}
	\pw GEp=\pw GpE\implies
	\pvat{V}{E}{T,p}=-\pvat{p_\electric}{p}{T,E}.
\end{equation}
对于各向同性、无自发极化的线性电介质，$p_\electric=\varepsilon_0 V\chi_\electric E$。弱场下上式右端正比于$E$，故恒温恒压下体积变化的首项正比于$E^2$。这种形变称为\textbf{电致伸缩}。

\paragraph{磁介质}

外磁场方向固定，以$m$记样品沿场方向的总磁矩，$B$记外场的磁感应强度。采用磁功$B\d m$的约定，有
\[
	\d U=T\d S-p\d V+B\d m.
\]
% 但应该注意的是，磁感应强度并不能作为平衡条件。
以外场为自变量的Gibbs自由能为$G=U-TS+pV-Bm$，其微分为
\[
	\d G=-S\d T+V\d p-m\d B
\]
Maxwell关系
\[
	\pw GBT=\pw GTB\implies
	\pvat{S}{B}{T,p}=\pvat{m}{T}{B,p}.
\]
在恒压下令$C_{p,B}=T(\p S/\p T)_{p,B}$。可逆绝热过程满足$\d S=0$，故
\[
	\pvat{T}{B}{S,p}=-\frac T{C_{p,B}}\pvat{m}{T}{B,p}.
\]
弱场下的Curie顺磁体满足$m\propto B/T$。固定正向外场时，升温使磁矩减小，故$(\p T/\p B)_{S,p}>0$，可逆绝热退磁使样品降温。这是\textbf{磁热效应}的一个实例。

